\documentclass{ximera}
\title{Rain and Einsten's Railway Carriage}

\newcommand{\pskip}{\vskip 0.1 in}

\begin{document}
\begin{abstract}
An introduction to relative motion.
\end{abstract}
\maketitle

\section{Einstein's Railway Carriage}


The purpose of mechanics is to desribe how bodies change their position in space with ``time''. I should load my conscience with grave sins against the sacred spirit of lucidity were I to formulate the aims of mechanics in this way, without serious reflection and detailed explanation. Let us proceed to disclose these sins.

 It is not clear to be understood here what is to be meant by ``position'' and ``space''. I stand at the window of a railway carriage which is travelling uniformly, and drop a stone on the embankment, without throwing it. Then, disregarding the influence of the air resistance, I see the stone descend in a straight line. A pedestrian who observes the misdeed from a footpath notices that the stone falls to earth in a parabolic curve. I now ask: Do the ``positions'' traversed by the stone lie in ``reality'' on a straight line or on a parabola? Moreover, what is meant here by motion ``in space''? From the considerations of the previous section, the answer is self-evident. In the first place, we entirely shun the vague word ``space'', of which, we must honestly acknowledge, we cannot form the slightest conception, and replace it by ``motion relative to a practically rigid body of reference''. The positions relative to the body of reference (railway carriage or embankment) have already been defined in detail in the previous section. If instead of ``body of reference'' we insert ``system of coordinates'', which is a useful idea for mathematical description, we are in a position to say: The stone traverses a straight line relative to a system of coordinates rigidly attached to the carriage, but relative to a system of coordinates rigidly attached to the ground (embankment) it describes a parabola. With the aid of this example it is clearly seen that there is no such thing as an independently existing trajectory (lit. ``path- curve''), but only a trajectory relative to a particular body of reference. 

From  Chapter 3, \emph{Space and Time in Classical Mechanics}, Albert Einstein's \emph{Relativity: The Special and General Theory}, 1916.


\begin{question} \label{Edsfdf4tDE}

Access Desmos interactive at
 
\href{https://www.desmos.com/calculator/ga6zcv9vkr}{Falling Rain 2}

 
\begin{onlineOnly}
    \begin{center}
\desmos{ga6zcv9vkr}{900}{600}
\end{center}
\end{onlineOnly}

\end{question}

\end{document}

